\documentclass[11pt]{article}

\usepackage{amsmath, amssymb, amsthm}
\usepackage{stmaryrd}   % \llbracket, \rrbracket, \fatsemi
\usepackage{mathtools}
\usepackage[margin=1in]{geometry}

% --- notation ---------------------------------------------------------------
\newcommand{\sem}[1]{\llbracket #1 \rrbracket}      % denotation
\newcommand{\Then}{\mathbin{\fatsemi}}              % THEN operator
\newcommand{\ThenG}[2]{\mathbin{\fatsemi}_{[#1,\,#2]}} % THEN with gap [gmin,gmax]
\newcommand{\Empty}{\bot}                           % empty / undefined result
\newcommand{\size}[1]{\lvert #1 \rvert}
\newcommand{\R}[5]{\langle #1,\, #2,\, #3,\, #4,\, #5 \rangle} % split repr.
\newcommand{\gmin}{g_{\min}}
\newcommand{\gmax}{g_{\max}}

\theoremstyle{definition}
\newtheorem{definition}{Definition}
\theoremstyle{plain}
\newtheorem{theorem}{Theorem}
\newtheorem{lemma}{Lemma}

\title{The \textsc{Then} Operation on Interval Sets\\
       \large (adjacency and bounded-gap concatenation)}
\author{}
\date{}

\begin{document}
\maketitle

\noindent
This note states, in human-readable form, the definitions of the \textsc{Then}
(sequential concatenation) operator---both the gap-free version and the
generalisation that allows a bounded gap between the two intervals---together
with the equivalences that are machine-checked in
\texttt{Intervalset/Then.lean} and \texttt{Intervalset/ThenWithGap.lean}.
Sections~\ref{sec:then-single}--\ref{sec:then-correct} treat the gap-free case;
Section~\ref{sec:gap} generalises to a bounded gap and recovers the gap-free
case as the specialisation $\gmin = \gmax = 0$.

% ===========================================================================
\section{Intervals and interval sets}

\begin{definition}[Interval]
An \emph{interval} $I = [a, b]$ is a pair of natural numbers with $a \le b$.
Its \emph{size} is $\size{I} = b - a + 1 \ge 1$.
\end{definition}

\begin{definition}[Interval set]
An \emph{interval set} is a finite set of intervals. We write membership as
$I \in A$.
\end{definition}

% ===========================================================================
\section{\textsc{Then} on single intervals}
\label{sec:then-single}

Two intervals concatenate exactly when the second begins on the frame
immediately after the first ends:
\[
  [a_1, b_1] \Then [a_2, b_2] \;=\;
  \begin{cases}
    [a_1, b_2] & \text{if } a_2 = b_1 + 1,\\[2pt]
    \Empty     & \text{otherwise.}
  \end{cases}
\]

% ===========================================================================
\section{\textsc{Then} on interval sets: enumeration semantics}

The reference (``ground truth'') semantics enumerates all pairs, one from each
set, and keeps the non-empty concatenations:
\[
  A \Then B \;=\;
  \bigl\{\, I_1 \Then I_2 \;\bigm|\; I_1 \in A,\; I_2 \in B,\;
            I_1 \Then I_2 \neq \Empty \,\bigr\}.
\]
This is correct but quadratic in $\size{A}\cdot\size{B}$.

% ===========================================================================
\section{The compact representation}

\begin{definition}[Split interval set]
A \emph{split interval set} $R = \R{s}{l}{r}{e}{c}$ denotes the set of all
intervals that start in $[s, l]$, end in $[r, e]$, and have size at least $c$:
\[
  \sem{R} \;=\;
  \bigl\{\, [a, b] \;\bigm|\;
            s \le a \le l,\;\;
            r \le b \le e,\;\;
            b - a + 1 \ge c \,\bigr\}.
\]
\end{definition}

\noindent
Geometrically, in the $(\text{start},\text{stop})$ plane $\sem{R}$ is the
axis-aligned box $[s, l] \times [r, e]$ with the lower-left corner clipped by
the diagonal half-plane $\text{stop} - \text{start} + 1 \ge c$.

% ===========================================================================
\section{\textsc{Then} on the compact representation}

Let $A = \R{s_1}{l_1}{r_1}{e_1}{c_1}$ and $B = \R{s_2}{l_2}{r_2}{e_2}{c_2}$,
with $c_1, c_2 \ge 1$.\footnote{Assuming $c \ge 1$ is no loss of generality:
every interval has size $\ge 1$, so $c = 0$ and $c = 1$ denote the same set.}

\paragraph{Feasibility.}
A result interval $[a,b]$ exists only if some split point $m$ (the stop of the
first half, with $m+1$ the start of the second) is simultaneously legal for
both operands. Such an $m$ can exist at all only when
\[
  \Phi(A, B)\;:\equiv\;
  \max(r_1,\, s_2 - 1)\;\le\;\min(e_1,\, l_2 - 1).
\]
If $\Phi$ fails, $A \Then B = \Empty$.

\paragraph{Result.}
When $\Phi(A,B)$ holds, the concatenation is again a single split interval set:
\[
  A \Then B \;=\; \R{\,s_1\,}{\,L\,}{\,R\,}{\,e_2\,}{\,c_1 + c_2\,},
\]
where
\[
  L \;=\; \min\!\bigl(l_1,\; e_1 - c_1 + 1,\; l_2 - c_1\bigr),
  \qquad
  R \;=\; \max\!\bigl(r_2,\; r_1 + c_2,\; s_2 + c_2 - 1\bigr).
\]

\noindent
Intuition for the bounds:
\begin{itemize}
  \item \textbf{Start} $a$: inherited lower bound $s_1$; the upper bound $L$ is
        the tightest of ``$A$'s own cap'' $l_1$, ``leave room for $c_1$ frames
        before $A$'s latest stop'' $e_1 - c_1 + 1$, and ``leave room for $c_1$
        frames before $B$'s latest start'' $l_2 - c_1$.
  \item \textbf{Stop} $b$: inherited upper bound $e_2$; the lower bound $R$ is
        the tightest of ``$B$'s own floor'' $r_2$, ``$c_2$ frames past $A$'s
        earliest stop'' $r_1 + c_2$, and ``$c_2$ frames past $B$'s earliest
        start'' $s_2 + c_2 - 1$.
  \item \textbf{Size}: the two halves contribute $c_1$ and $c_2$ frames with no
        gap and no overlap, so the result needs size $\ge c_1 + c_2$
        (i.e.\ $g_{\min} = 0$ in the gap-free case).
\end{itemize}

% ===========================================================================
\section{Correctness}
\label{sec:then-correct}

\begin{theorem}[\texttt{thenOp\_spec}]
For all split interval sets $A, B$ with $c_1, c_2 \ge 1$ and every interval
$I$,
\[
  I \in \sem{A \Then B}
  \quad\Longleftrightarrow\quad
  \exists\, I_1 \in \sem{A},\, I_2 \in \sem{B}.\;\; I_1 \Then I_2 = I.
\]
Equivalently, the compact operation computes exactly the enumeration semantics:
\[
  \sem{A \Then B} \;=\; \sem{A} \Then \sem{B}.
\]
\end{theorem}

\noindent
The proof is fully formalised in Lean~4 / Mathlib and depends only on the
standard axioms (\texttt{propext}, \texttt{Classical.choice},
\texttt{Quot.sound}); see \texttt{thenOp\_spec} in
\texttt{lean-prove/Intervalset/Then.lean}.

\paragraph{Implementation caveat (natural-number arithmetic).}
Over $\mathbb{N}$ the subtractions in $L$ must be evaluated as
$e_1 + 1 - c_1$ and $l_2 - c_1$ (not $e_1 - c_1 + 1$), and the feasibility
guard must additionally check $c_1 \le l_2$, so that truncated subtraction is
never silently floored to $0$. Both are reflected in the verified definition.

% ===========================================================================
\section{\textsc{Then} with a bounded gap}
\label{sec:gap}

We now generalise to a \emph{gap} of $g$ frames between the two intervals,
where $g$ is constrained to a window $[\gmin, \gmax]$ (with $\gmin \le \gmax$).
The gap-free operator is the special case $\gmin = \gmax = 0$.

% ---------------------------------------------------------------------------
\subsection{Single intervals}

Writing the gap as $g = a_2 - b_1 - 1$ (the number of frames strictly between
the two intervals), the concatenation keeps the pair iff $g \in [\gmin, \gmax]$:
\[
  [a_1, b_1] \ThenG{\gmin}{\gmax} [a_2, b_2] \;=\;
  \begin{cases}
    [a_1, b_2] & \text{if } b_1 + 1 + \gmin \le a_2 \le b_1 + 1 + \gmax,\\[2pt]
    \Empty     & \text{otherwise.}
  \end{cases}
\]
The result interval $[a_1, b_2]$ \emph{spans the gap}, so its size is
$\size{[a_1,b_1]} + \size{[a_2,b_2]} + g$. The enumeration semantics on sets is
again the set of all non-empty pairwise concatenations,
$A \ThenG{\gmin}{\gmax} B = \{\, I_1 \ThenG{\gmin}{\gmax} I_2 \mid I_1 \in A,\,
I_2 \in B,\, \neq \Empty \,\}$.

% ---------------------------------------------------------------------------
\subsection{Compact representation}

Let $A = \R{s_1}{l_1}{r_1}{e_1}{c_1}$, $B = \R{s_2}{l_2}{r_2}{e_2}{c_2}$ with
$c_1, c_2 \ge 1$ and $\gmin \le \gmax$.

\paragraph{Feasibility.}
A result now needs a first stop $m$ \emph{and} a gap $g \in [\gmin, \gmax]$ with
the second interval starting at $m + 1 + g$. Both must be jointly realisable:
\[
  \Phi_g(A, B)\;:\equiv\;
  \underbrace{r_1 \le e_1}_{\text{first stop exists}}
  \;\wedge\;
  \max(r_1 + \gmin,\; s_2 - 1) \;\le\; \min(e_1 + \gmax,\; l_2 - 1).
\]
If $\Phi_g$ fails, $A \ThenG{\gmin}{\gmax} B = \Empty$.

\paragraph{Result.}
When $\Phi_g(A,B)$ holds,
\[
  A \ThenG{\gmin}{\gmax} B \;=\;
  \R{\,s_1\,}{\,L\,}{\,R\,}{\,e_2\,}{\,c_1 + c_2 + \gmin\,},
\]
where
\[
  L \;=\; \min\!\bigl(l_1,\; e_1 + 1 - c_1,\; l_2 - c_1 - \gmin\bigr),
  \qquad
  R \;=\; \max\!\bigl(r_2,\; r_1 + c_2 + \gmin,\; s_2 + c_2 - 1\bigr).
\]

\noindent
Changes relative to the gap-free case:
\begin{itemize}
  \item \textbf{Size}: the gap contributes at least $\gmin$ extra frames, so the
        result needs size $\ge c_1 + c_2 + \gmin$. This is exactly the
        $g_{\min}$ term: $c' = c_1 + c_2 + \gmin$.
  \item \textbf{Start} $a$: the term $e_1 + 1 - c_1$ is \emph{unchanged}---a
        larger gap does not relax it, because the first stop $m$ is capped by
        $e_1$ no matter how big the gap is. Only the $B$-facing term tightens to
        $l_2 - c_1 - \gmin$ (room must be left for the gap as well).
  \item \textbf{Stop} $b$: the cross term becomes $r_1 + c_2 + \gmin$ (the gap
        pushes $B$'s earliest reachable stop further to the right).
\end{itemize}

\paragraph{Derivation sketch.}
Membership of $[a,b]$ requires a first stop $m \in [\max(r_1, a + c_1 - 1),\,
e_1]$ and a gap $g \in [\gmin, \gmax]$, and every constraint coming from $B$
depends on $m$ and $g$ only through the sum $t = m + g$. Since $m$ and $g$ range
over intervals independently, $t$ ranges over
$[\max(r_1, a{+}c_1{-}1) + \gmin,\; e_1 + \gmax]$, and a witness exists iff this
range is non-empty and meets the $B$-required range
$[s_2 - 1,\; \min(l_2 - 1,\; b - c_2)]$. Expanding that condition and projecting
onto $(a,b)$ yields $\Phi_g$, $L$, $R$ and the size bound above.

% ---------------------------------------------------------------------------
\subsection{Correctness}

\begin{theorem}[\texttt{thenGapOp\_spec}]
For all split interval sets $A, B$ with $c_1, c_2 \ge 1$, all $\gmin \le \gmax$,
and every interval $I$,
\[
  I \in \sem{A \ThenG{\gmin}{\gmax} B}
  \;\Longleftrightarrow\;
  \exists\, I_1 \in \sem{A},\, I_2 \in \sem{B}.\;\;
  I_1 \ThenG{\gmin}{\gmax} I_2 = I.
\]
\end{theorem}

\noindent
As before the proof is fully formalised in Lean~4 / Mathlib and uses only the
standard axioms (\texttt{propext}, \texttt{Classical.choice},
\texttt{Quot.sound}); see \texttt{thenGapOp\_spec} in
\texttt{lean-prove/Intervalset/ThenWithGap.lean}.

\begin{theorem}[\texttt{thenGapOp\_zero\_zero}: conservative extension]
For all $A, B$,
\[
  A \ThenG{0}{0} B \;=\; A \Then B .
\]
\end{theorem}

\noindent
On the compact representation this holds \emph{definitionally} (the Lean proof
is \texttt{rfl} and depends on \emph{no} axioms): with $\gmin = \gmax = 0$ every
$+\,0$ and $-\,0$ disappears and the two operations are syntactically the same
after reduction. Hence the bounded-gap algebra is a conservative extension of
the gap-free one, and the two correctness theorems necessarily agree at
$\gmin = \gmax = 0$.

\paragraph{Implementation caveat (natural-number arithmetic).}
As in the gap-free case the $L$ subtractions are evaluated as $e_1 + 1 - c_1$
and $l_2 - c_1 - \gmin$, and the feasibility guard additionally checks
$c_1 + \gmin \le l_2$ so the term $l_2 - c_1 - \gmin$ is never silently floored
to $0$ in $\mathbb{N}$.

\end{document}
